\documentclass[11pt,a4paper]{article}
\usepackage[utf8]{inputenc}
\usepackage[british]{babel}
\usepackage{geometry}
\geometry{margin=1in}
\usepackage{graphicx}
\usepackage{booktabs}
\usepackage{longtable}
\usepackage{float}
\usepackage{amsmath}
\usepackage{amssymb}
\usepackage{hyperref}
\usepackage{caption}
\usepackage{array}
\usepackage{ragged2e}

\hypersetup{colorlinks=true, linkcolor=black, citecolor=black, urlcolor=blue}

\title{Risk Management \\ GWP3\\[0.3em]}
\author{Ge Zhang \\ Alberto Hernandez-Espinosa}
\date{\today}

\begin{document}

\maketitle



\vspace{0.6em}

\section*{On AI usage}

This research made use of artificial intelligence tools for: (i) literature search and the 
verification of academic sources and datasets; (ii) to
refine and optimise algorithmic implementations and validation; and (iii) for grammar, 
spelling and style corrections. All intellectual contributions,
theoretical framing, empirical analysis and final interpretations are entirely our own.

% \tableofcontents
%\newpage

%==========================================================================
\section{Introduction}

This report closes the mini-capstone begun in GWP1 and GWP2. The underlying study is still Danish A. 
Alvi's \textit{Application of Probabilistic Graphical Models in
Forecasting Crude Oil Price} (University College London, 2018). That claims that the price
of crude oil can be predicted by (i) a discrete panel of macroeconomic and physical market time
series into latent regimes with hidden Markov models, (ii) learning a belief network over those
discretised variables with a hill-climbing search, and (iii) performing probabilistic inference on
the network to obtain the regime of the oil price in the coming month.

Forecasting the price of crude oil is a well studied problem, and the literature shows that its
moves come from different demand and supply shocks (Kilian; Baumeister and Kilian).

We introduce our findings and conclusions at the end of this report.

%==========================================================================
\section{Training, validation and testing sets}

Our dataset contains 199 month-end observations from January 2010 to July 2026, 
organized as follows:
\begin{itemize}
    \item \textbf{Variables}: Seven time series (see Figure~\ref{fig:panel_gauss}), 
    including front-month WTI and 
    Brent crude futures, and the US Dollar Index (all sourced from Yahoo Finance), 
    plus the Consumer Price Index, effective federal funds rate, industrial 
    production index, and WTI spot price (all sourced from the Federal 
    Reserve Economic Data service; see United States, Federal Reserve Bank of St. Louis).
    \item \textbf{Temporal alignment}: All series were standardised to the 
    last business day of each month to ensure consistent time alignment across variables.
    \item \textbf{Outlier handling}: Extreme values were capped using
    the interquartile-range rule with a multiplier of three to limit the influence of outliers.
    \item \textbf{Gap filling}: Any missing observations in the aligned 
    series were imputed using time-based linear interpolation.
\end{itemize}

The dataset is divided chronologically into three windows. Table~\ref{tab:roles} states this 
division which is also explained in the following sections.

\begin{table}[H]
\centering
\caption{The division of the dataset between the three windows.}
\label{tab:roles}
\renewcommand{\arraystretch}{1.2}
\begin{tabular}{@{}lp{5.1cm}p{4.6cm}r@{}}
\toprule
\textbf{Set} & \textbf{What is estimated on it} & \textbf{How often it may be consulted} &
\textbf{Months} \\
\midrule
Training   & Hidden Markov parameters, network structure, conditional probability tables &
Freely & 139 \\
Validation & Nothing & Freely, but only to choose between specifications already fitted & 30 \\
Testing    & Nothing & Once, after every choice has been frozen & 30 \\
\bottomrule
\end{tabular}
\end{table}

As Table~\ref{tab:roles} shows we created a chronological division of the dataset into 
three windows in such a way that training precedes validation and validation precedes testing.


\subsection{The training set}

The training set is the only window where anything is estimated. Three estimation problems are
solved on it, those summarized in Table~\ref{tab:training_tasks}.

\begin{table}[H]
\centering
\caption{Estimation tasks performed on the training set}
\label{tab:training_tasks}
\renewcommand{\arraystretch}{1.3}
\begin{tabular}{@{}p{3.5cm}p{\dimexpr\textwidth-3.5cm-2\tabcolsep\relax}@{}}
\toprule
\textbf{Estimation Task} & \textbf{Implementation Details} \\
\midrule
Three-state hidden Markov model parameters $(\Pi, A, B)$ (one per time series) & Fitted via the Baum-Welch algorithm (Rabiner); re-estimated from 10 random initializations to avoid local optima, retaining the solution with the highest log-likelihood. \\
\midrule
Belief network structure & Searched via hill climbing over the discretised training dataset (Chickering). \\
\midrule
Conditional probability tables (for each node given its parents) & Estimated under a K2 prior (Cooper and Herskovits) to avoid zero probability assignments for unobserved (but not impossible) state configurations, which would occur with maximum-likelihood estimation given the small sample size. \\
\bottomrule
\end{tabular}
\end{table}

The window, as is shown in Table~\ref{tab:roles}, runs from January 2010 to July 2021, 139 
months or 69.8 per cent of the sample, and
contains the two structural sections that define the modern crude market: the collapse 
of late 2014, driven by the expansion of US shale supply, and the demand destruction of 2020. 
A training window that does not contain both sections would produce a model with no 
experience of crisis dynamics.


\subsection{The validation set}

The validation set estimates nothing; it decides between specifications that were all estimated
on the training set. It is needed because a specification chosen only by how well it fits the
training data is chosen partly for how well it reproduces noise. Here we made 4 decisions:

\begin{enumerate}
\item the scoring function of the structure search, among the K2 score, the Bayesian Dirichlet
      equivalent uniform score (Heckerman, Geiger and Chickering) and the discrete Bayesian
      information criterion;
\item the maximum in-degree permitted in the search are 2, 3, unrestricted, which
      supports how many parents a node may acquire and then how quickly the conditional
      probability tables become sparse;
\item whether the search is seeded with an expert graph encoding standard energy-market
      economics, or started from an empty graph;
\item the decision threshold that converts the posterior distribution over regimes into a
      position.
\end{enumerate}

We used the validation set in the same way as the source paper, which experiments with ``different
Bayesian Dirichlet scoring functions such as the BDeu or K2 and the Bayesian Information
Criterion'' (Alvi 41) and states that ``the validation step is to adjust the model if 
the error is too high'' (Alvi 62). We used a window  that runs from August 2021 
to January 2024, thirty months covering the post-pandemic recovery and the energy 
shock that followed the invasion of Ukraine. We compared 18
configurations on it for each of the two specifications reported in
Section~\ref{sec:results}.

\subsection{The testing set}

The testing set measures, and only measures. We used it in this way:

\begin{enumerate}
\item it is opened once, after the discretisation, the structure, the parameters and the decision
      threshold have all been frozen;
\item nothing is adjusted along the way, because each new look would turn it into a second
      validation set;
\item the error it reports is the number quoted, without further qualification or selection among
      variants.
\end{enumerate}

The window runs from February 2024 to July 2026, 30 months, following and mirroring our target
reference, that suggest that at least thirty observations are needed before a
proportion can be estimated with any confidence (Alvi 41). Then we adopt a 70:15:15 division 
rather than the 80:10:10 as in the source reference, because with 199 months, an 80:10:10 
division would leave twenty months on each side, and the exact binomial confidence 
interval around an accuracy of, say, 60 per cent would then span roughly 36 to 81 per cent, 
which is too wide to support any conclusion at all. 
Then, our division set as 70:15:15 preserves the paper's logic while giving
30 months to each hold-out window; even so, as Section~\ref{sec:results} shows, the intervals
remain wide.

%==========================================================================
\section{Allocation of the data and comparison of the hold-out sets}

Table~\ref{tab:split} shows the basic statistics of each window, and
Figure~\ref{fig:allocation} shows the division against the price series.

\begin{table}[H]
\centering
\caption{Allocation of the 199 monthly observations and descriptive statistics of the WTI spot
price in each window.}
\label{tab:split}
\begin{tabular}{@{}lrlrrrrr@{}}
\toprule
\textbf{Set} & \textbf{Months} & \textbf{Period} & \textbf{Share} & \textbf{Mean} &
\textbf{Mean} & \textbf{Vol.} & \textbf{Range} \\
 & & & (\%) & price & log ret. & log ret. & (USD) \\
\midrule
Training   & 139 & 2010-01 to 2021-07 & 69.8 & 69.34 & \phantom{-}0.0001 & 0.1231 & 19.2--113.4 \\
Validation &  30 & 2021-08 to 2024-01 & 15.1 & 83.61 & \phantom{-}0.0037 & 0.0950 & 66.1--114.4 \\
Testing    &  30 & 2024-02 to 2026-07 & 15.1 & 73.39 & -0.0000 & 0.1196 & 57.3--108.6 \\
\bottomrule
\end{tabular}
\end{table}

\begin{figure}[H]
\centering
\includegraphics[width=\textwidth]{figures/data_allocation.png}
\caption{Chronological allocation of the sample. The division respects the arrow of time: every
observation used for model selection postdates every observation used for estimation, and every
observation used for testing postdates both.}
\label{fig:allocation}
\end{figure}

\subsection{Are the two hold-out windows comparable?}

The question here is if the validation and testing windows describe the market in a reasonable way.
 We test this on the monthly log returns of the WTI spot
price with three tests:
\begin{itemize}
    \item the two-sample Kolmogorov-Smirnov test for equality of the full distributions,
    \item Levene's test for equality of variances between the two samples,
    \item and Welch's test for equality of the sample means.
\end{itemize}
Table~\ref{tab:ks} reports the results.

\begin{table}[H]
\centering
\caption{Tests of distributional equality between the windows, monthly log returns of the WTI spot
price.}
\label{tab:ks}
\begin{tabular}{@{}lrrrr@{}}
\toprule
\textbf{Comparison} & \textbf{KS statistic} & \textbf{KS $p$} & \textbf{Levene $p$} &
\textbf{Welch $p$} \\
\midrule
Training vs validation & 0.144 & 0.647 & 0.921 & 0.860 \\
Training vs testing    & 0.123 & 0.812 & 0.972 & 0.996 \\
Validation vs testing  & 0.172 & 0.791 & 0.953 & 0.895 \\
\bottomrule
\end{tabular}
\end{table}

None of the tests rejects at the five per cent level. This means that the three
windows do not look statistically different in their central tendency, their spread, or their overall
shape. That is good in specific way, because it suggests that choosing the model on the 
validation window is not the same as tuning it on a market that is completely different 
from the one used for final testing. Even so, two qualifications are important. First, 
failing to reject is not the same as proving that the windows are equal. With only 
thirty observations, these tests have limited power, so they would detect only a 
fairly large difference. Second, similarity in the distribution of returns is not 
the same as similarity in the underlying economic environment. The validation
window reflects a war-related supply issue, whereas the testing window reflects a phase of
gradual normalisation. Those two situations may generate return distributions 
that look similar on paper even though they come from different economic forces. 
This distinction matters because a belief network is meant to learn the causal 
structure behind the data, not just match the unconditional distribution. For 
that reason, the comparison supports the validation procedure, but it does not 
prove that a structure learned in one regime will necessarily carry over to another.

\subsection{How the regimes are distributed within each window}

The distributional tests compare the raw returns, but what the model actually 
forecasts is the regime, so the two windows must also be compared on their regime 
 composition. Table~\ref{tab:shares} reports it for the discretisation retained 
 in Section~\ref{sec:regimes}.

\begin{table}[H]
\centering
\caption{Composition of each window by decoded regime of the WTI spot price, log-return emissions
(per cent of months).}
\label{tab:shares}
\begin{tabular}{@{}lrrrr@{}}
\toprule
\textbf{Set} & \textbf{Months} & \textbf{Bear} & \textbf{Stagnant} & \textbf{Bull} \\
\midrule
Training   & 138 & 21.7 & 65.9 & 12.3 \\
Validation &  30 & 23.3 & 63.3 & 13.3 \\
Testing    &  30 & 13.3 & 73.3 & 13.3 \\
\bottomrule
\end{tabular}
\end{table}

The training and validation windows are closely similar, but the testing window is softer than
any other: the stagnant state is 73.3 per cent of its months against 63.3 per cent in the
validation window, and the bear state is 13.3 per cent against 23.3 per cent. The effects
are summarised below and discussed in Section~\ref{sec:results}.

These findings are summarised in Table~\ref{tab:window_implications}, which formalises the key differences between the validation and testing windows that affect model evaluation:

\begin{table}[H]
\centering
\caption{Key distinctions between validation and testing windows that impact model performance
assessment. The baselines are measured on the window itself; the 72.4 per cent reported in
Table~\ref{tab:accuracy} is the same rule measured over the 29 months of the one-month-ahead
target.}
\label{tab:window_implications}
\renewcommand{\arraystretch}{1.3}
\begin{tabular}{@{}p{3.5cm}p{5.2cm}p{5.2cm}@{}}
\toprule
\textbf{Metric} & \textbf{Validation Window} & \textbf{Testing Window} \\
\midrule
Always-stagnant baseline accuracy & 63.3\% & 73.3\% \\
Number of directional (non-stagnant) months & 11 (36.7\% of window) & 8 (26.7\% of window) \\
Total economically informative months for forecasting & 11 & 8 \\
\midrule
Core evaluation implication & Model accuracy must outperform the lower 63.3\% baseline to demonstrate skill & The higher 73.3\% baseline means raw accuracy can appear inflated without additional context; small sample of directional months produces wide confidence intervals that cannot be resolved by model refinement alone \\
\bottomrule
\end{tabular}
\end{table}

\subsection{Comparison by purpose}

Beyond statistics, the two sets differ in what a poor result on each would mean:
 
\begin{itemize}
    \item A poor result on the validation set is informative and actionable: it tells 
    us to change the scoring function, the in-degree limit or the discretisation, 
    and we may do so as often as we wish.
    \item A poor result on the testing set is terminal for the specification: acting 
    on it would corrupt the estimate.
\end{itemize}

This
asymmetry why we use both, and it is why the 18 configurations reported in
Section~\ref{sec:replication} were compared on the validation window alone.

%==========================================================================
\section{Regime detection: from prices to discrete evidence}
\label{sec:regimes}

A discrete belief network needs inputs that take a fixed and finite number of values, whereas our
panel is made of continuous series. Following the source paper, we converted each
series into a sequence of three regimes: bear, stagnant, and bull. This was done with a
three-state hidden Markov model estimated on the training window (Rabiner). The states were then ordered by
the arithmetic mean of the underlying monthly change, so that the labels keeps an economic
interpretation rather than being arbitrary state numbers.

Two emission schemes are estimated and compared.

\paragraph{Parity emissions (the source paper's scheme).} The series is first differenced, and each
change is reduced to its sign. This creates a two-symbol alphabet: one symbol for an increase and
zero for a decrease (Alvi 42). The resulting model is a discrete-emission hidden Markov model
with three hidden states.

\paragraph{Log-return emissions (our modification).} In our modified version, the emission is the
monthly log return itself. To limit the influence of extreme observations, returns are capped at
three training-sample standard deviations and then modelled with a Gaussian density with diagonal
covariance. The transition matrix is given a sticky Dirichlet prior, implemented by adding
pseudo-counts of five to the diagonal (Fox et al.). In plain terms, this encodes the prior belief that market
regimes usually persist for several months rather than switching every month.

In both cases, decoding is performed \textit{in real time}. That means the regime assigned to month
$t$ is obtained by running the Viterbi recursion using only the emissions observed up to and
including month $t$. This is a change from the source paper, which
decodes each hold-out window as a whole (Alvi 60, 63). When an entire window is decoded at once,
the state assigned to a given month is influenced in part by what happens later in the sample. The
regime series then contains information that would not actually have been available at the time,
and any forecast built on that series is contaminated by look-ahead bias. In our first
implementation we reproduced that choice and obtained a nominal one-month-ahead accuracy of
100 per cent, which is a sign of leakage and not of a good model.

\subsection{What the two schemes produce}

Table~\ref{tab:hmm} summarises the estimated models for the WTI spot price, and
Figure~\ref{fig:transmat} displays the two transition matrices.

\begin{table}[H]
\centering
\caption{Hidden Markov models of the WTI spot price estimated on the training window.}
\label{tab:hmm}
\begin{tabular}{@{}lrr@{}}
\toprule
 & \textbf{Parity emissions} & \textbf{Log-return emissions} \\
\midrule
Log-likelihood & $-92.37$ & $138.53$ \\
Average diagonal of the transition matrix & 0.000 & 0.742 \\
Regime switches over 198 months & 189 & 52 \\
Mean monthly log return, bear state & $-$ & $-0.129$ \\
Mean monthly log return, stagnant state & $-$ & $+0.014$ \\
Mean monthly log return, bull state & $-$ & $+0.152$ \\
\bottomrule
\end{tabular}
\end{table}

\begin{figure}[H]
\centering
\includegraphics[width=\textwidth]{figures/transition_matrices.png}
\caption{Estimated transition matrices for the WTI spot price. Under parity emissions the diagonal
is empty: the chain is obliged to leave its state at every step. Under log-return emissions the
mass concentrates on the diagonal, which is the behaviour a market regime is supposed to exhibit.}
\label{fig:transmat}
\end{figure}

The main conclusion is direct. The parity specification does not recover meaningful
market regimes: because it sees only whether price moved up or down, it produces a sequence that
switches state in 189 of 198 months and behaves almost like an alternation. By contrast, the
log-return specification recovers persistent and economically interpretable states, with average
persistence of 0.74, about four months of expected duration, and a much lower number of switches
(52 over the full sample), which is closer to the behaviour described in the regime-switching
literature (Hamilton; Ang and Timmermann). For that reason, the log-return model is the specification that carries
forward into the analysis. Table~\ref{tab:regime_summary} summarises the comparison,
Table~\ref{tab:latent} gives the latent meaning of each state, Figure~\ref{fig:regimes}
shows the decoded sequence against the price, and Figure~\ref{fig:regimes_parity} shows the same
series under the parity scheme.

\begin{table}[H]
\centering
\caption{Why the log-return specification is more interpretable than the parity specification}
\label{tab:regime_summary}
{\small\setlength{\tabcolsep}{4pt}\renewcommand{\arraystretch}{1.15}\begin{tabular}{>{\raggedright\arraybackslash}p{2.8cm} >{\raggedright\arraybackslash}p{5.0cm} >{\raggedright\arraybackslash}p{5.0cm}}
\toprule
Feature & Parity emissions & Log-return emissions \\
\midrule
What the model sees & Only whether the monthly change is positive or negative. & The size and sign of the monthly log return, after capping extremes. \\
Decoded behaviour & Regimes switch in 189 of 198 months, so states alternate almost every month. & Regimes switch 52 times, which is much more consistent with persistent market phases. \\
Economic meaning & Weak: the states mainly separate signs, not genuine market conditions. & Strong: the states map naturally into bear, stagnant, and bull phases. \\
Persistence & Zero probability of remaining in the same state. & Average persistence is 0.74, implying a regime duration of about four months. \\
Overall interpretation & The model does not recover usable market regimes. & The model produces regimes an analyst can interpret and use. \\
\bottomrule
\end{tabular}}
\end{table}




\begin{table}[H]
\centering
\caption{Latent meaning of the hidden states under log-return emissions, training window.}
\label{tab:latent}
\begin{tabular}{@{}lrrrr@{}}
\toprule
\textbf{State} & \textbf{Months} & \textbf{Mean log return} & \textbf{Volatility} &
\textbf{Share (\%)} \\
\midrule
Bear     & 30 & $-0.1289$ & 0.1463 & 21.7 \\
Stagnant & 91 & $+0.0143$ & 0.0532 & 65.9 \\
Bull     & 17 & $+0.1518$ & 0.1350 & 12.3 \\
\bottomrule
\end{tabular}
\end{table}

\begin{figure}[H]
\centering
\includegraphics[width=\textwidth]{figures/regimes_full_gaussian.png}
\caption{Decoded regimes of the WTI spot price under log-return emissions. The bear state captures
the second half of 2014, the first months of 2020 and the 2025 correction; the bull state captures
the 2021 recovery and the first quarter of 2022.}
\label{fig:regimes}
\end{figure}

\begin{figure}[H]
\centering
\includegraphics[width=\textwidth]{figures/regimes_full_parity.png}
\caption{Decoded regimes of the same series under the parity emissions of the dissertation. The
colour changes almost every month; no episode of the past sixteen years is identified as a
sustained bear or bull market.}
\label{fig:regimes_parity}
\end{figure}

\subsection{The seven series, not only the price}

The belief network uses the discretised version of all seven variables, so
Figures~\ref{fig:panel_gauss} and~\ref{fig:panel_parity} show every series coloured by the regime
its own hidden Markov model attributes to it. The labels are relative to each series, so for a
trending macroeconomic index ``bear'' means the state of slowest growth rather than a falling
market.

\begin{figure}[H]
\centering
\includegraphics[width=\textwidth]{figures/regimes_panel_gaussian.png}
\caption{Decoded regimes of all seven modelled series under log-return emissions. The two crude
futures and the spot price show blocks lasting several months; the Consumer Price Index, the
industrial production index and the dollar index are almost monochromatic.}
\label{fig:panel_gauss}
\end{figure}

\begin{figure}[H]
\centering
\includegraphics[width=\textwidth]{figures/regimes_panel_parity.png}
\caption{The same seven series under the parity emissions of the dissertation. The dollar index
changes regime in 197 of 198 months and the spot price in 189; Brent, by contrast, collapses onto a
single state. Neither behaviour is a market regime.}
\label{fig:panel_parity}
\end{figure}

The panel view reveals something that the single-series plot in Figure~\ref{fig:regimes} does not
show clearly. Under log-return emissions, the three oil price series behave in a way that broadly
matches the idea of a regime process: WTI futures switch regime 46 times, Brent 35 times, and the
spot price 52 times over 198 months. The macroeconomic variables behave very differently. The
Consumer Price Index never leaves its first state, industrial production spends 97 per cent of the
sample in one state, and the dollar index spends 88 per cent in one state. At monthly frequency
these series are close to monotonic, so once they are modelled with Gaussian emissions and a prior
that favours persistence, one state absorbs almost every observation. After discretisation, the
variable is therefore close to constant.

That matters because a nearly constant variable contributes almost no information to the belief
network. This is why the structure search in Section~\ref{sec:results} drops the Consumer Price
Index, the federal funds rate, and the dollar index, while retaining the price-based nodes. In
other words, the macroeconomic part of the model is muted before the network is even estimated.
The limiting step is the discretisation scheme, not the belief network itself. A more promising
approach for future work would be to discretise trending macroeconomic series relative to a moving
growth benchmark rather than using the raw monthly change.

Under parity emissions the problem appears in the opposite form, but it is just as damaging. The
dollar index alternates in 197 of 198 months, while Brent, whose differenced monthly changes are
almost always of the same sign, collapses into a single state and switches only twice in sixteen
years. One series changes state almost every month; the other barely changes state at all. Neither
pattern matches a useful economic notion of regime.

That is the main economic message of the two figures. A regime that lasts only one month is not
really a regime, and a system that reports a new market state every month cannot guide a position
intended to be held for a quarter. The advantage of the improved specification is therefore not
just statistical neatness. It changes the kind of state sequence that the downstream network is
being asked to learn and predict.

%==========================================================================
\section{Re-running the Bayesian network with hill climbing}
\label{sec:replication}

We reproduced the original paper's protocol as closely as possible. A node called
\texttt{forecast}, which is a copy of the discretised WTI spot-price regime, is added to
the training data. The network structure is then learned by hill climbing, the conditional
probability tables are estimated under a K2 prior, and the model is asked for the most probable
state of \texttt{forecast} given the seven remaining variables. The resulting sequence of
inferences is then shifted forward by one month so that it can be read as a prediction. This is the
same one-month shift that the original paper applies in its own implementation (Alvi 62--63). The
computations use open-source implementations of hidden Markov models and belief networks (Ankan
and Panda).

We compared 18 configurations on the validation window, combining three scoring functions,
three in-degree limits, and the presence or absence of the expert seed. The configuration that was
retained used the K2 score, unrestricted in-degree, and the expert graph as a seed. It produced a
network with 24 edges, in which the Markov blanket of the forecast node (Koller and Friedman)
includes all seven remaining variables. Figure~\ref{fig:network_k2} shows the network learned by the unseeded K2 search on the
training window.

\begin{figure}[H]
\centering
\includegraphics[width=0.85\textwidth]{figures/network_k2_train.png}
\caption{Belief network over the discretised variables learned by hill climbing with the K2 score
on the training window. The forecast node, shaded, is attached to the futures and macroeconomic
nodes, reproducing the densely connected topology reported by the dissertation (Alvi 60).}
\label{fig:network_k2}
\end{figure}

\subsection{Were the results of the paper replicated?}

Table~\ref{tab:replication} places our figures beside those reported in the dissertation.

\begin{table}[H]
\centering
\caption{Replication of the validation and testing errors of Alvi (2018, 62--63).}
\label{tab:replication}
\begin{tabular}{@{}llrrr@{}}
\toprule
\textbf{Study} & \textbf{Set} & \textbf{Months} & \textbf{Error (\%)} & \textbf{Accuracy (\%)} \\
\midrule
Alvi (2018)      & Validation & 28 & 67.86 & 32.14 \\
Alvi (2018)      & Testing    & 28 & 42.86 & 57.14 \\
This replication & Validation & 30 & 86.67 & 13.33 \\
This replication & Testing    & 30 & 76.67 & 23.33 \\
\bottomrule
\end{tabular}
\end{table}

The short answer is: \textbf{partially}.

\textbf{What does replicate is the procedure and its qualitative outcome}. The full pipeline runs on
an independent dataset, the hill-climbing search produces a densely connected graph of the same
general kind as the one reported in the original paper, the forecast node acquires the futures and
macroeconomic variables as parents, and the validation error is very high, just as it is in the
paper, where roughly two-thirds of the validation months are misclassified.

\textbf{What does not replicate is the size of the testing error}. The original paper reports 
42.9 per cent
error, or 57.1 per cent accuracy, and presents that result as encouraging. Our replication
produces 76.7 per cent error, or 23.3 per cent accuracy. That is \textit{worse} than
the 33.3 per cent accuracy an uninformed guess would achieve if the three states were equally
likely. When a model performs consistently below chance, the problem is not noise. It
is a sign that the model is systematically wrong, and that systematic error needs an explanation.

We attribute this underperformance to the interaction between the parity discretisation 
(Section~\ref{sec:regimes}) and the one-month shifting convention. 
The network only produces a \textit{nowcast}: because the \texttt{forecast} node 
is an exact copy of the current month's target regime, the structure search 
learns this deterministic relationship and replicates the current regime 
nearly perfectly. Shifting this perfect nowcast forward by one month converts 
its accuracy into a measure of regime persistence: \textit{it only correctly predicts 
next month's regime if the regime does not change}. 

The parity emission framework 
produces regimes with near-zero persistence by construction (the estimated chain 
switches states almost every month), so the protocol never actually measures the 
network's predictive power, it only measures the persistence of its discretised labels. 
The original paper's 57 per cent testing accuracy is plausibly 
explained by its sample producing more persistent regimes from the same 
procedure. \textbf{In our implementation we cannot confirm that the network 
genuinely forecasts future states}.

This leads to one of our conclusions: the forecast has to come from the network itself, and not
from the persistence of the labels.

%==========================================================================
\section{Results: the accuracy of the crude oil forecast}
\label{sec:results}

Three specifications were used to evaluate the network's predictive power on identical 
months of the testing window.

\begin{description}
\item[Replication model.] This is the original protocol: parity emissions, a
duplicated forecast node, and the inferred state shifted in advance by one month.

\item[Improved model.] We used as target the regime of the \textit{next} month, so the network is
asked to predict rather than reconstruct. The evidence is the current regime of the
seven series under log-return emissions. Model selection is again carried out on the validation
window. The retained specification uses the discrete Bayesian information criterion, a maximum
in-degree of two, and no expert seed, which produces a sparse four-edge network
(Figure~\ref{fig:network_pred}), \texttt{WTI\_Spot}\,$\rightarrow$\,\texttt{WTI\_CL}\,$\rightarrow$\,\texttt{Brent\_BZ},
\texttt{WTI\_CL}\,$\rightarrow$\,\texttt{forecast} and
\texttt{Brent\_BZ}\,$\rightarrow$\,\texttt{Ind\_Prod}, in which the Markov blanket of the forecast
node reduces to the front-month WTI futures contract.

\item[Expected-return rule.] The posterior distribution over next month's regime, produced by the
improved model, is converted into a trading position by maximising the expected monthly log return
implied by that posterior, using the regime means in Table~\ref{tab:latent}. A threshold of one
per cent, selected on the validation window, keeps the rule flat when the expected move is small.
\end{description}

\begin{figure}[H]
\centering
\includegraphics[width=0.85\textwidth]{figures/network_predictive.png}
\caption{The improved predictive network. Selection on the validation window favours a sparse
graph: with 138 training months and three states per node, dense networks exhaust the data in
their conditional probability tables.}
\label{fig:network_pred}
\end{figure}

\subsection{Classification accuracy}

Table~\ref{tab:accuracy} reports the main results, and Figure~\ref{fig:bars} compares them with
three benchmarks: an uninformed guess, always predicting the modal regime from the training
window, and simply repeating last month's regime.

\begin{table}[H]
\centering
\caption{One-month-ahead forecast of the WTI spot regime. Balanced accuracy is the mean of the
per-regime recalls and is not inflated by the dominance of the stagnant state.}
\label{tab:accuracy}
\begin{tabular}{@{}llrrrr@{}}
\toprule
\textbf{Model} & \textbf{Set} & \textbf{Months} & \textbf{Accuracy} & \textbf{Balanced} &
\textbf{Macro} \\
 & & & (\%) & \textbf{accuracy (\%)} & \textbf{F1 (\%)} \\
\midrule
Replication          & Validation & 30 & 13.33 & 17.32 & 13.03 \\
Replication          & Testing    & 30 & 23.33 & 22.35 & 19.26 \\
Improved             & Validation & 29 & 58.62 & 34.39 & 30.95 \\
Improved             & Testing    & 29 & 75.86 & 41.67 & 39.31 \\
Expected-return rule & Testing    & 29 & 68.97 & 45.24 & 45.46 \\
\midrule
Uninformed guess     & Testing    & 29 & 34.48 & --- & --- \\
Majority regime      & Testing    & 29 & 72.41 & --- & --- \\
Persistence          & Testing    & 29 & 79.31 & --- & --- \\
\bottomrule
\end{tabular}
\end{table}

\begin{figure}[H]
\centering
\includegraphics[width=\textwidth]{figures/accuracy_comparison.png}
\caption{Accuracy on the testing window. The dashed line is the accuracy of an uninformed guess
over three states. Neither belief-network specification beats the persistence benchmark.}
\label{fig:bars}
\end{figure}

Three points stand out, and the third is the most important.

First, \textbf{the replication model performs worse than chance}. Its exact 95 per cent binomial confidence
interval is 9.9 to 42.3 per cent, for the structural reason established in
Section~\ref{sec:replication}.

Second, \textbf{the improved model classifies 75.9 per cent of months correctly, with an exact interval of
56.5 to 89.7 per cent}. This is above the majority benchmark of 72.4 per cent, but still below the
persistence benchmark of 79.3 per cent. Exact McNemar tests on the paired month-by-month outcomes
(McNemar) give $p = 1.00$ against both benchmarks. In a sample this short, the belief network is therefore
statistically indistinguishable from the very simple rule of repeating last month's regime.

Third, \textbf{the balanced accuracy of the improved model is only 41.7 per cent}.
Figure~\ref{fig:conf} explains why. Under the maximum a posteriori rule, the model answers
``stagnant'' in twenty-six of twenty-nine months. It correctly identifies the 21 stagnant months,
but it captures only one of the four bear months and none of the four bull months.
Figure~\ref{fig:timeline} shows the same pattern month by month: the model follows the long
stagnant stretch of 2024 and early 2025, and misses the turning points of mid-2025 and 2026.
A rule that is
correct three times out of four because it keeps saying ``nothing much will happen'' may look good
on a headline accuracy measure, but it is not useful in practice. The distinction between raw
accuracy and decision usefulness is exactly the distinction a risk manager needs to keep in view.

\begin{figure}[H]
\centering
\begin{minipage}{0.49\textwidth}
\includegraphics[width=\textwidth]{figures/confusion_modelA_test.png}
\end{minipage}
\begin{minipage}{0.49\textwidth}
\includegraphics[width=\textwidth]{figures/confusion_modelB_test.png}
\end{minipage}
\caption{Confusion matrices on the testing window: replication model (left) and improved model
(right). The improved model concentrates its answers on the stagnant regime; the replication model
distributes its errors across all three.}
\label{fig:conf}
\end{figure}

\begin{figure}[H]
\centering
\includegraphics[width=\textwidth]{figures/timeline_modelB_test.png}
\caption{Month-by-month forecasts of the improved model on the testing window. Correct months are
marked with a circle and incorrect months with a cross. The model tracks the long stagnant stretch
of 2024 and early 2025 and misses the turning points of mid-2025 and 2026.}
\label{fig:timeline}
\end{figure}

\subsection{From probabilities to decisions}

The weakness of the maximum a posteriori rule is a decision problem, not a probability problem.
The posterior itself, shown in Figure~\ref{fig:posterior}, does move in sensible ways: the
probability of the bear state rises before the mid-2025 correction, and the probability of the
bull state rises before the 2026 rebound, even though neither state ever becomes the most likely
one. If we keep only the modal regime, we discard that information. The reason is that the loss
function behind the maximum a posteriori rule treats all three kinds of mistakes as equally
costly, whereas in a commodity book the cost of missing a twelve per cent decline is plainly not
the same as the cost of missing a one per cent drift.

The expected-return rule restores the economic interpretation. Writing $p_t(k)$ for the posterior
probability of regime $k$ in month $t$ and $\mu_k$ for the mean monthly log return of that regime
estimated on the training window, the rule takes a long position when $\sum_k p_t(k)\mu_k$
exceeds a threshold $\tau$, a short position when it falls below $-\tau$, and no position
otherwise. The threshold is set to $\tau = 0.01$ on the validation window.

The result, shown in Table~\ref{tab:accuracy}, is that raw accuracy falls from 75.9 to
69.0 per cent, while balanced accuracy rises from 41.7 to 45.2 per cent and macro F1 rises from
39.3 to 45.5 per cent. The rule makes a directional call in nine of the twenty-nine months, which
is 31 per cent coverage, and it is correct in 66.7 per cent of those calls.
Figure~\ref{fig:strategy} traces the value of one barrel traded under those signals. Starting from
79.2 US dollars, it ends the testing window at 157.0 dollars, compared with 70.6 dollars for a
barrel held throughout.

\begin{figure}[H]
\centering
\includegraphics[width=\textwidth]{figures/posterior_modelC.png}
\caption{Posterior probability of each regime for the coming month over the testing window, with
the realised regime marked above. The bear probability rises before the mid-2025 decline without
ever becoming modal, which is what the maximum a posteriori rule discards.}
\label{fig:posterior}
\end{figure}

\begin{figure}[H]
\centering
\includegraphics[width=\textwidth]{figures/strategy_modelC.png}
\caption{Value of one barrel traded on the expected-return rule against a barrel held throughout
the testing window. Transaction costs, financing and slippage are excluded and the path rests on
nine decisions; the exhibit illustrates the decision rule and is not a performance claim.}
\label{fig:strategy}
\end{figure}

That backtest still has to be read carefully. It excludes transaction costs, financing, and
slippage; it assumes execution at the month-end price; it is based on only nine directional
decisions, so a single reversed call would materially change the ending value; and the threshold
was selected on a thirty-month validation window, which itself introduces selection bias. What the
exhibit does show legitimately is that the same posterior can produce materially different, and
more useful, behaviour when it is evaluated through an economic loss function rather than a
symmetric classification rule. In that sense, the choice of loss function deserves at least as
much attention as the choice of scoring metric for the network.

\subsection{Summary of the empirical findings}

\begin{enumerate}
\item The dissertation's original protocol does not transfer successfully to an independent
sample. In our data it delivers 23.3 per cent accuracy, which is below chance, because it combines
non-persistent regimes with a shifted nowcast.
\item Redefining regimes by the size of the monthly move, adding a prior that favours
persistence, and decoding strictly causally produces economically interpretable states and raises
raw accuracy to 75.9 per cent.
\item Even so, that accuracy is statistically indistinguishable from a persistence rule
($p = 1.00$), and the balanced accuracy of 41.7 per cent shows that the model defaults too often
to the ``stagnant'' state.
\item Converting the posterior into a trading position through an expected-return rule raises
balanced accuracy to 45.2 per cent and yields directional calls that are correct two times out of
three, at the cost of some raw accuracy.
\item Over a sample of only thirty months, none of these differences is statistically significant.
The honest conclusion is that the belief network may be extracting a small amount of information
about crude-oil direction, but the sample is too short to establish that the information is real.
\end{enumerate}

%==========================================================================
\section{Assessment of the contributions claimed by the dissertation}

The dissertation closes with eight claimed contributions (Alvi 66--67). We assess each against the
evidence the document itself provides and against what we observed in reproducing the work. Our
verdicts are summarised in Table~\ref{tab:contrib} and argued below.

\begin{table}[H]
\centering
\caption{Assessment of the eight claimed contributions of Alvi (2018).}
\label{tab:contrib}
\renewcommand{\arraystretch}{1.25}
\begin{tabular}{@{}p{0.5cm}p{7.4cm}p{2.6cm}p{4.2cm}@{}}
\toprule
\textbf{\#} & \textbf{Claim} & \textbf{Verdict} & \textbf{Principal evidence} \\
\midrule
1 & Bayesian views replace EGARCH-M views in a Black-Litterman model & Not achieved &
pp. 66--67; no allocation implemented \\
2 & Time series discretised by hidden Markov models as inputs to belief networks & Achieved &
Secs. 3.3, 4.2--4.2.5, pp. 42--43, 54--58 \\
3 & A working trading mechanism capable of independent decisions & Partially achieved &
Sec. 4.3.1, pp. 64--65 \\
4 & An autonomous system requiring no prior expert knowledge & Partially achieved &
Sec. 4.2.6, p. 59 \\
5 & The abstraction of the Python libraries clarifies the design process & Achieved &
Secs. 2.2.1--2.2.2, pp. 28--29 \\
6 & A global macro strategy returning more than high-frequency or fixed-income funds & Not achieved
& Sec. 4.3.1, p. 65 \\
7 & Better models of energy markets to inform policy & Partially achieved &
Sec. 3.4, pp. 43--45; figure p. 60 \\
8 & Amalgamation of existing research to increase alpha for commodity traders & Partially achieved
& Ch. 5, p. 67 \\
\bottomrule
\end{tabular}
\end{table}

\paragraph{1. Replacing EGARCH-M derived views with Bayesian model derived views in the
Black-Litterman framework (p. 66).} The claim refers to Beach and Orlov, whose 2007 study derives
Black-Litterman views from an EGARCH-M model. The dissertation proposes the substitution in its
conclusions and cites the source, but no Black-Litterman allocation appears anywhere in the
document: there is no prior equilibrium return vector, no view matrix, no confidence matrix and no
resulting portfolio. Chapter 4 ends with a single-asset long/short rule on one barrel of crude
(pp. 64--65), which is not a portfolio allocation. We judge this contribution not achieved; it is
a research proposal rather than a result. Had it been carried out it would have been the most
valuable of the eight, because a probabilistic model that outputs a full posterior is a natural
source of the view confidences that Black-Litterman requires and that practitioners usually set by
hand.

\paragraph{2. Using time series discretised by hidden Markov models as inputs to belief networks
(p. 66).} This is the methodological core of the work and it is genuinely implemented: Section 3.3
sets out the design (pp. 42--43), Sections 4.2 to 4.2.5 estimate the models and decode the regimes
(pp. 54--58), the regime-switch plot of the spot price on p. 56 shows the result, and the
discretised frames feed the structure search. We reproduced the mechanism
successfully, which is the strongest possible evidence that it works. We judge it achieved, with
one qualification about novelty: regime-switching representations of commodity markets long
predate the dissertation, and the specific contribution is the coupling of the decoded states to a
structure-learned belief network rather than the discretisation itself. This matters because it is
the one component of the paper that a practitioner can adopt directly, and our own results show
that the value of the coupling depends entirely on how the emissions are defined.

\paragraph{3. A working trading mechanism capable of independent decision making (p. 66).}
Section 4.3.1 implements a rule that goes long on a forecast bull regime, short on a forecast bear
regime and stays flat otherwise, and plots the resulting value of one barrel against the price and
against the EIA Short Term Energy Outlook (figure on p. 65; United States, Energy Information
Administration). The mechanism therefore exists and it does
make decisions without human intervention. It is not, however, a trading system in any operational
sense: transaction costs, financing, slippage and position limits are absent, the evaluation covers
twenty-eight months, and no risk-adjusted measure is computed. Our own equivalent exhibit
(Figure~\ref{fig:strategy}) suffers the same limitations, which is why we present it as an
illustration. Partially achieved.

\paragraph{4. An autonomous system requiring no prior expert knowledge beyond dataset selection
(p. 66).} The dissertation's own code contradicts the claim in its strong form. On p. 59 the hill
climbing is initialised from an expert graph derived from EIA relationships --- six edges linking
OPEC and non-OPEC production and OECD and non-OECD demand to the spot price --- and the search then
attaches the remaining variables. A search seeded with an expert structure is not free of expert
knowledge; the seed determines which local optimum the greedy search reaches. Our experiments make
this concrete: seeded and unseeded searches returned different networks under the same score, and
on the improved specification the unseeded sparse network was the one selected on the validation
window. The claim holds for the unseeded variant, which the dissertation does not evaluate
separately. Partially achieved.

\paragraph{5. The level of abstraction of the Python modelling libraries clarifies the design
process (p. 66--67).} Sections 2.2.1 and 2.2.2 (pp. 28--29) do explain how the libraries express
the theory, and the exposition is clear and pedagogically useful. We accept the claim as achieved
while noting that it is a presentational contribution about third-party software rather than a
scientific one, and that it has aged: the library used for the hidden Markov models is no longer
maintained, which forced both our Group Work Project 2 and this project to substitute a different
implementation. The episode is a reminder that reproducibility claims resting on a specific
software stack decay quickly.

\paragraph{6. A systematic, event-driven global macro strategy capable of generating higher returns
than funds based on high-frequency or fixed-income strategies (p. 67).} No comparison with any fund,
index or strategy of either type appears in the document. The only benchmark shown is the EIA Short
Term Energy Outlook (p. 65), which is a forecast, not an investable strategy, and the comparison
covers a handful of episodes in 2016 and 2017 discussed qualitatively. The author's own
concluding sentence on that page concedes that ``further research is necessary to statistically
conclude with confidence if the model performs better than the EIA's STEO''. We judge this claim not
achieved, and we regard it as the weakest of the eight: it compares asset classes and strategy
families on the strength of a twenty-eight-month single-asset backtest with no risk adjustment.

\paragraph{7. Better models of energy markets, allowing policy makers to understand the underlying
structure (p. 67).} Section 3.4 (pp. 43--45) and the learned graph on p. 60 do deliver a structure
that can be inspected and discussed, which is more than a black-box forecaster offers, and this is
the genuine appeal of the graphical approach for a policy audience. The qualification is
stability. A structure intended to inform policy must be robust, and ours was not: across
eighteen configurations the number of edges ranged from two to twenty-five, and the parents of the
forecast node changed with the scoring function. Greedy search over a super-exponential space on
139 observations cannot be expected to recover a stable causal graph, and the dissertation does not
report any stability analysis such as bootstrap edge frequencies. Partially achieved; the framework
is right, the evidential standard is not yet met.

\paragraph{8. Amalgamating research from several disciplines and applying it to the commodity
markets to increase alpha for quantitative traders (p. 67).} The synthesis is real and is the
document's principal strength: probabilistic graphical models, regime detection, macroeconomic
data engineering and a trading application are brought into a single coherent pipeline, and the
dissertation is a useful map of that territory. The clause about increasing alpha is unsupported,
since no risk-adjusted return, benchmark comparison or significance test is reported anywhere.
Partially achieved.

\subsection{Why the achieved contributions matter}

Two of the eight claims survive scrutiny intact, and both are worth having. The coupling of hidden
Markov discretisation with structure-learned belief networks (claim 2) matters because it gives an
energy desk something that neither a GARCH model nor a neural network provides: a forecast that
arrives with an explicit, inspectable statement of which variables it depended on. When the model
says the coming month is likely to be bearish, the Markov blanket says why, and a risk committee
can argue with the reasoning rather than with the number. The expository contribution (claim 5)
matters for a narrower but practical reason: it lowers the cost of entry for practitioners who need
to audit such a model rather than build one.

The four partially achieved claims share a single deficiency, and it is not conceptual but
evidential. The trading mechanism, the autonomy of the structure search, the policy relevance of
the graph and the alpha claim would all be settled by evidence the dissertation could have
produced and did not: transaction costs, a seeded-versus-unseeded comparison, bootstrap edge
stability, and a risk-adjusted benchmark. Our own results reinforce the point from the other
direction. We produced some of that evidence --- exact confidence intervals, McNemar tests against
persistence and majority benchmarks, a coverage-and-hit-rate decomposition --- and it showed that
what looked like a respectable 75.9 per cent accuracy is statistically indistinguishable from
repeating last month's answer. The lesson is that in this problem the evidential apparatus is not a
formality appended to the modelling; it is what separates a result from an artefact.

%==========================================================================
\section{Discussion}

\textit{This section is written for a reader who takes investment decisions and is deliberately
free of technical vocabulary.}

\subsection{What the study did and what it found}

Using information that anyone can obtain free of charge --- the price of oil, the strength of the
dollar, inflation, the interest rate set by the Federal Reserve and a measure of industrial
activity --- we built a system that describes the oil market each month as being in one of three
conditions: falling, flat or rising. The system then estimates the chance that each condition will
prevail in the month ahead.

Tested on the two and a half years the system had never seen, it named the prevailing condition
correctly in roughly three months out of four. That figure sounds impressive and is largely
misleading, because the market was flat in most of those months and a rule that always said
``flat'' would have scored almost as well. The more honest measure is how the system behaved at the
moments that mattered: when it took a clear directional view, which it did in about one month in
three, it was right two times out of three, and acting on those views over the period would have
been considerably more profitable than simply holding a barrel of oil.

We also repeated the original study on which this work is based, and found that its method does not
transfer to our data. The reason is worth stating plainly, because it is the most useful thing we
learned: the original method defined the market's condition purely by whether the price went up or
down last month. Defined that way, the condition changes almost every month, and a description that
changes every month cannot support a decision that is held for a quarter. Once the condition was
defined by the size of the move rather than merely its direction, the descriptions became stable,
recognisable and usable --- the system independently identified the 2014 collapse, the 2020 crash,
the 2021 recovery and the 2025 correction as distinct episodes.

\subsection{The recommended course of action}

We do not recommend trading on this system alone, and we do recommend using it in three specific
ways.

\textbf{Use it to size risk, not to pick trades.} Its most reliable output is not a prediction but
a warning: the small number of months in which the evidence points clearly up or down. In those
months, an energy book should be positioned more decisively; in the remaining two-thirds of months,
the sensible action is to reduce activity rather than to invent a view.

\textbf{Use it as a challenge to the house view.} Because the system states which indicators drove
its conclusion, a committee can interrogate it. When it disagrees with the desk's opinion, the
disagreement identifies precisely which piece of evidence the two readings weigh differently, and
that conversation is valuable even when the system turns out to be wrong.

\textbf{Keep any position modest and reviewed monthly.} The evidence base is thirty months and nine
directional calls. The apparent profitability excludes trading costs. On that evidence, positions
should be small, should carry a predefined loss limit, and should be reconsidered as each month's
figures are published.

\subsection{The factors that move the result}

Three factors dominate, and a decision maker should watch them in this order.

\textbf{The definition of the market's condition.} This is the single largest influence on the
answer --- larger than any refinement of the model that consumes it. A definition based on the size
of price moves produces conditions that persist for about four months; a definition based on
direction alone produces noise. Anyone commissioning a similar system should scrutinise this
choice first.

\textbf{The dollar and the cost of money.} Oil is priced in dollars, so a stronger dollar makes it
more expensive everywhere else and dampens demand; higher policy rates raise the cost of carrying
inventory and of financing production. Together they push the market towards the falling condition,
and they were among the indicators the system attached to its forecast.

\textbf{Industrial activity.} It stands in for demand. When it weakens, the flat condition tends to
give way to the falling one.

A fourth factor is the reason for caution: events with no statistical precedent. Wars, production
decisions by exporting countries and the storage crisis of 2020 dominate the price for months at a
time and are not in the information the system reads. It will always be late to these, and no
amount of statistical refinement will change that. It should be used as one considered opinion
among several, and never as the only one.

%==========================================================================
\section{Conclusion and practical takeaways}

The mini-capstone set out to determine whether probabilistic graphical models can forecast the
price of crude oil more effectively than conventional time-series methods. Our answer, on sixteen
years of monthly data and a thirty-month hold-out window, is a qualified no with a useful
remainder.

\begin{enumerate}
\item \textbf{The framework is sound and its published implementation is fragile.} We could not
replicate the dissertation's testing accuracy, and we identified why: its regime definition
produces non-persistent labels, and its evaluation protocol converts an accurate nowcast into a
systematically wrong forecast. Practitioners adopting this literature should test the persistence
of their regime labels before anything else.
\item \textbf{Discretisation dominates model choice.} Changing the emission from the sign of the
monthly move to its magnitude, with a prior favouring persistence, moved accuracy from 23.3 to
75.9 per cent. No change of scoring function, in-degree limit or prior in the belief network came
close to that effect.
\item \textbf{Accuracy is the wrong headline.} An accuracy of 75.9 per cent that is
indistinguishable from persistence ($p = 1.00$) and that carries a balanced accuracy of 41.7 per
cent describes a model that defaults to inaction. Balanced accuracy, coverage and hit rate should
be reported alongside accuracy in any regime-forecasting study.
\item \textbf{The loss function is part of the model.} Reading the same posterior with an
expected-return rule rather than a maximum a posteriori rule raised balanced accuracy by three and
a half points, produced directional calls that were right two times in three, and changed the
economic outcome materially. This is the cheapest available improvement in the whole pipeline.
\item \textbf{Sample size is the binding constraint.} Monthly macroeconomic data give roughly two
hundred observations per decade and a half. Thirty-month hold-out windows produce confidence
intervals thirty points wide. Any future work in this direction should either move to a higher
frequency, pool across related commodities, or adopt a rolling-origin evaluation that reuses the
sample honestly rather than a single hold-out split (Hyndman and Athanasopoulos).
\end{enumerate}

For a risk manager, the practical value of the exercise lies less in the forecast than in its
transparency. A belief network states which variables its conclusion rests upon; when it is wrong,
the reason is inspectable and the model can be corrected rather than merely retrained. In a market
whose structure changes every few years, that property may matter more than a few points of
accuracy.

%==========================================================================
\newpage
\section*{Works Cited}
\addcontentsline{toc}{section}{Works Cited}

\noindent
Alvi, Danish A. \textit{Application of Probabilistic Graphical Models in Forecasting Crude Oil
Price}. 2018. University College London, dissertation. \textit{arXiv},
\url{https://arxiv.org/abs/1804.10869}.

\vspace{0.6em}\noindent
Ang, Andrew, and Allan Timmermann. ``Regime Changes and Financial Markets.'' \textit{Annual Review
of Financial Economics}, vol. 4, no. 1, 2012, pp. 313--337.

\vspace{0.6em}\noindent
Ankan, Ankur, and Abinash Panda. ``pgmpy: Probabilistic Graphical Models Using Python.''
\textit{Proceedings of the 14th Python in Science Conference}, 2015, pp. 6--11.

\vspace{0.6em}\noindent
Baumeister, Christiane, and Lutz Kilian. ``Forecasting the Real Price of Oil in a Changing World: A
Forecast Combination Approach.'' \textit{Journal of Business \& Economic Statistics}, vol. 33,
no. 3, 2015, pp. 338--351.

\vspace{0.6em}\noindent
Beach, Steven L., and Alexei G. Orlov. ``An Application of the Black--Litterman Model with
EGARCH-M-Derived Views for International Portfolio Management.'' \textit{Financial Markets and
Portfolio Management}, vol. 21, no. 2, 2007, pp. 147--166.

\vspace{0.6em}\noindent
Chickering, David Maxwell. ``Optimal Structure Identification with Greedy Search.'' \textit{Journal
of Machine Learning Research}, vol. 3, 2002, pp. 507--554.

\vspace{0.6em}\noindent
Cooper, Gregory F., and Edward Herskovits. ``A Bayesian Method for the Induction of Probabilistic
Networks from Data.'' \textit{Machine Learning}, vol. 9, no. 4, 1992, pp. 309--347.

\vspace{0.6em}\noindent
Fox, Emily B., et al. ``A Sticky HDP-HMM with Application to Speaker Diarization.'' \textit{The
Annals of Applied Statistics}, vol. 5, no. 2A, 2011, pp. 1020--1056.

\vspace{0.6em}\noindent
Hamilton, James D. ``A New Approach to the Economic Analysis of Nonstationary Time Series and the
Business Cycle.'' \textit{Econometrica}, vol. 57, no. 2, 1989, pp. 357--384.

\vspace{0.6em}\noindent
Heckerman, David, Dan Geiger, and David M. Chickering. ``Learning Bayesian Networks: The
Combination of Knowledge and Statistical Data.'' \textit{Machine Learning}, vol. 20, no. 3, 1995,
pp. 197--243.

\vspace{0.6em}\noindent
Hyndman, Rob J., and George Athanasopoulos. \textit{Forecasting: Principles and Practice}. 3rd ed.,
OTexts, 2021.

\vspace{0.6em}\noindent
Kilian, Lutz. ``Not All Oil Price Shocks Are Alike: Disentangling Demand and Supply Shocks in the
Crude Oil Market.'' \textit{American Economic Review}, vol. 99, no. 3, 2009, pp. 1053--1069.

\vspace{0.6em}\noindent
Koller, Daphne, and Nir Friedman. \textit{Probabilistic Graphical Models: Principles and
Techniques}. MIT Press, 2009.

\vspace{0.6em}\noindent
McNemar, Quinn. ``Note on the Sampling Error of the Difference between Correlated Proportions or
Percentages.'' \textit{Psychometrika}, vol. 12, no. 2, 1947, pp. 153--157.

\vspace{0.6em}\noindent
Rabiner, Lawrence R. ``A Tutorial on Hidden Markov Models and Selected Applications in Speech
Recognition.'' \textit{Proceedings of the IEEE}, vol. 77, no. 2, 1989, pp. 257--286.

\vspace{0.6em}\noindent
United States, Energy Information Administration. \textit{Short-Term Energy Outlook}. U.S.
Department of Energy, 2026, \url{https://www.eia.gov/outlooks/steo/}.

\vspace{0.6em}\noindent
United States, Federal Reserve Bank of St. Louis. \textit{Federal Reserve Economic Data (FRED)}.
2026, \url{https://fred.stlouisfed.org/}.

\end{document}
